\documentclass[../../main.tex]{subfiles}
\begin{document}

\begin{center}
  {\fontsize{48pt}{58pt}\selectfont \textbf{2002}\normalsize\textbf{年度}}
\end{center}

\begin{table}[htb]
  \centering
  \caption{出題テーマと難易度}
  \begin{tabular}{|c|p{9cm}|c|}
    \hline
    \textbf{問題番号} & \textbf{テーマ} & \textbf{難易度} \\
    \hline
    第1問 & 札を箱に振り分ける確率 & \\
    \hline
    第2問 & 楕円と円が4点で交わる条件 & \\
    \hline
    第3問 & 四面体の射影と共円条件 & \\
    \hline
    第4問 & 有理数での値から多項式の係数を決める & \\
    \hline
    第5問 & ペル方程式を保つ数列と極限 & \\
    \hline
    第6問 & 積分方程式で定まる関数 & \\
    \hline
  \end{tabular}
\end{table}

\begin{table}[htb]
  \centering
  \caption{解答}
  \renewcommand{\arraystretch}{1.4}
  \begin{tabular}{|c|c|p{9cm}|}
    \hline
    \textbf{問題番号} & \textbf{小問} & \textbf{答え} \\
    \hline
    第1問 & -- & $\dfrac{1}{2(n-2)!}$ \\
    \hline
    第2問 & -- & 範囲は解答の図を参照（境界を除く） \\
    \hline
    第3問 & -- & 証明問題 \\
    \hline
    第4問 & -- & 証明問題 \\
    \hline
    第5問 & (1) & $a_n^2-2b_n^2=1$ \\
    \cline{2-3}
    & (2) & $\sqrt{2}$ \\
    \hline
    第6問 & -- & $f(x)=x+1-\dfrac{4\pi+8}{4+\pi^2}\sin x+\dfrac{8-4\pi}{4+\pi^2}\cos x$ \\
    \hline
  \end{tabular}
\end{table}

\paragraph{試験形態}
6問．試験時間は要確認．

\end{document}
